\begin{thebibliography}{99}
\footnotesize\itemsep=0pt

\bibitem{AlNi}
Alexeev V., Nikulin V.V., Del {P}ezzo and {$K3$} surfaces, \textit{MSJ
 Memoirs}, Vol.~15, \href{https://doi.org/10.1142/e002}{Mathematical Society of Japan}, Tokyo, 2006.

\bibitem{AlHa}
Altmann K., Hausen J., Polyhedral divisors and algebraic torus actions,
 \href{https://doi.org/10.1007/s00208-005-0705-8}{\textit{Math. Ann.}} \textbf{334} (2006), 557--607, \href{https://arxiv.org/abs/math.AG/0306285}{arXiv:math.AG/0306285}.

\bibitem{ArBrHaWr}
Arzhantsev I., Braun L., Hausen J., Wrobel M., Log terminal singularities,
 platonic tuples and iteration of {C}ox rings, \href{https://doi.org/10.1007/s40879-017-0179-8}{\textit{Eur.~J.~Math.}}
 \textbf{4} (2018), 242--312, \href{https://arxiv.org/abs/1703.03627}{arXiv:1703.03627}.

\bibitem{ArDeHaLa}
Arzhantsev I., Derenthal U., Hausen J., Laface A., Cox rings, \textit{Cambridge
 Stud. Adv. Math.}, Vol.~144, \href{https://doi.org/10.1017/CBO9781139175852}{Cambridge University Press}, Cambridge, 2015.

\bibitem{ArHaHeLi}
Arzhantsev I., Hausen J., Herppich E., Liendo A., The automorphism group of a
 variety with torus action of complexity one, \href{https://doi.org/10.17323/1609-4514-2014-14-3-429-471}{\textit{Mosc. Math.~J.}}
 \textbf{14} (2014), 429--471, \href{https://arxiv.org/abs/1202.4568}{arXiv:1202.4568}.

\bibitem{Ba}
B\"auerle A., Sharp degree bounds for fake weighted projective spaces,
 \href{https://arxiv.org/abs/2207.01709}{arXiv:2207.01709}.

\bibitem{BaHa}
B\"auerle A., Hausen J., Gorenstein, $\mathbb Q$-factorial, log terminal,
 non-toric {F}ano 3-folds of {P}icard number~1 with 2-torus action~-- a data
 file, {Z}enodo, available at \url{https://doi.org/10.5281/zenodo.7298788}.

\bibitem{BeHaHuNi}
Bechtold B., Hausen J., Huggenberger E., Nicolussi M., On terminal {F}ano
 3-folds with 2-torus action, \href{https://doi.org/10.1093/imrn/rnv190}{\textit{Int. Math. Res. Not.}} \textbf{2016}
 (2016), 1563--1602, \href{https://arxiv.org/abs/1412.8153}{arXiv:1412.8153}.

\bibitem{BrHae}
Braun L., H\"attig D., Canonical threefold singularities with a torus action of
 complexity one and {$k$}-empty polytopes, \href{https://doi.org/10.1216/rmj.2020.50.881}{\textit{Rocky Mountain~J.~Math.}}
 \textbf{50} (2020), 881--939, \href{https://arxiv.org/abs/1807.08022}{arXiv:1807.08022}.

\bibitem{BrKa}
Brown G., Kasprzyk A., Kawamata boundedness for {F}ano threefolds and the
 {G}raded {R}ing {D}atabase, \href{https://arxiv.org/abs/2201.07178}{arXiv:2201.07178}.

\bibitem{Cox}
Cox D.A., The homogeneous coordinate ring of a toric variety,
 \textit{J.~Algebraic Geom.} \textbf{4} (1995), 17--50,
 \href{https://arxiv.org/abs/alg-geom/9210008}{arXiv:alg-geom/9210008}.

\bibitem{CoLiSc}
Cox D.A., Little J.B., Schenck H.K., Toric varieties, \textit{Grad. Stud.
 Math.}, Vol.~124, \href{https://doi.org/10.1090/gsm/124}{Amer. Math. Soc.}, Providence, RI, 2011.

\bibitem{Dan}
Danilov V.I., The geometry of toric varieties, \href{https://dx.doi.org/10.1070/RM1978v033n02ABEH002305}{\textit{Russian Math. Surveys}}
 \textbf{33} (1978), 97--154.

\bibitem{Ful}
Fulton W., Introduction to toric varieties. {T}he {W}illiam {H}.~{R}oever
 lectures in geometry, \textit{Ann. of Math. Stud.}, Vol.~131, \href{https://doi.org/10.1515/9781400882526}{Princeton
 University Press}, Princeton, NJ, 1993.

\bibitem{HaHaHaSp}
H\"attig D., Hafner B., Hausen J., Springer J., Del {P}ezzo surfaces of
 {P}icard number one admitting a torus action, \href{https://arxiv.org/abs/2207.14790}{arXiv:2207.14790}.

\bibitem{HaHe}
Hausen J., Herppich E., Factorially graded rings of complexity one, in Torsors,
 \'etale Homotopy and applications to Rational Points, \textit{London Math.
 Soc. Lecture Note Ser.}, Vol.~405, \href{https://doi.org/10.1017/CBO9781139525350.011}{Cambridge Univ. Press}, Cambridge, 2013,
 414--428, \href{https://arxiv.org/abs/1005.4194}{arXiv:1005.4194}.

\bibitem{HaHiWr}
Hausen J., Hische C., Wrobel M., On torus actions of higher complexity,
 \href{https://doi.org/10.1017/fms.2019.35}{\textit{Forum Math. Sigma}} \textbf{7} (2019), e38, 81~pages,
 \href{https://arxiv.org/abs/1802.00417}{arXiv:1802.00417}.

\bibitem{HaKe}
Hausen J., Keicher S., A software package for {M}ori dream spaces,
 \href{https://doi.org/10.1112/S1461157015000212}{\textit{LMS~J.~Comput. Math.}} \textbf{18} (2015), 647--659,
 \href{https://arxiv.org/abs/1409.6929}{arXiv:1409.6929}.

\bibitem{HaSu}
Hausen J., S\"u{\ss} H., The {C}ox ring of an algebraic variety with torus
 action, \href{https://doi.org/10.1016/j.aim.2010.03.010}{\textit{Adv. Math.}} \textbf{225} (2010), 977--1012,
 \href{https://arxiv.org/abs/0903.4789}{arXiv:0903.4789}.

\bibitem{HiWa}
Hidaka F., Watanabe K., Normal {G}orenstein surfaces with ample anti-canonical
 divisor, \href{https://doi.org/10.3836/tjm/1270215157}{\textit{Tokyo~J.~Math.}} \textbf{4} (1981), 319--330.

\bibitem{HiWr}
Hische C., Wrobel M., On the anticanonical complex, \href{https://arxiv.org/abs/1808.01997}{arXiv:1808.01997}.

\bibitem{Il}
Ilten N., Mishna M., Trainor C., Classifying {F}ano complexity-one
 {$T$}-varieties via divisorial polytopes, \href{https://doi.org/10.1007/s00229-018-1036-x}{\textit{Manuscripta Math.}}
 \textbf{158} (2019), 463--486, \href{https://arxiv.org/abs/1710.04146}{arXiv:1710.04146}.

\bibitem{Is1}
Iskovskih V.A., Fano threefolds.~{I}, \href{https://doi.org/10.1070/IM1977v011n03ABEH001733}{\textit{Izv. Math.}} \textbf{11} (1977),
 485--527.

\bibitem{Is2}
Iskovskih V.A., Fano threefolds.~{II}, \href{https://doi.org/10.1070/IM1978v012n03ABEH001994}{\textit{Izv. Math.}} \textbf{12} (1978),
 469--506.

\bibitem{Ka}
Kasprzyk A.M., Canonical toric {F}ano threefolds, \href{https://doi.org/10.4153/CJM-2010-070-3}{\textit{Canad.~J.~Math.}}
 \textbf{62} (2010), 1293--1309, \href{https://arxiv.org/abs/0806.2604}{arXiv:0806.2604}.

\bibitem{MoMu}
Mori S., Mukai S., Classification of {F}ano {$3$}-folds with {$B_{2} \geq 2$},
 \href{https://doi.org/10.1007/BF01170131}{\textit{Manuscripta Math.}} \textbf{36} (1981), 147--162.

\bibitem{Nak}
Nakayama N., Classification of log del {P}ezzo surfaces of index two,
 \textit{J.~Math. Sci. Univ. Tokyo} \textbf{14} (2007), 293--498.

\bibitem{Pr1}
Prokhorov Yu.G., The degree of {$\mathbb Q$}-{F}ano threefolds, \href{https://doi.org/10.1070/SM2007v198n11ABEH003901}{\textit{Sb.
 Math.}} \textbf{198} (2007), 1683--1702.

\bibitem{Pr2}
Prokhorov Yu.G., {$\mathbb Q$}-{F}ano threefolds of large {F}ano index,~{I},
 \textit{Doc. Math.} \textbf{15} (2010), 843--872, \href{https://arxiv.org/abs/0812.1695}{arXiv:0812.1695}.

\bibitem{Pr3}
Prokhorov Yu.G., On {F}ano threefolds of large {F}ano index and large degree,
 \href{https://doi.org/10.1070/SM2013v204n03ABEH004304}{\textit{Sb. Math.}} \textbf{204} (2013), 4347--382.

\bibitem{Pr4}
Prokhorov Yu.G., {$\mathbb{Q}$}-{F}ano threefolds of index~7, \href{https://doi.org/10.1134/S0081543816060092}{\textit{Proc.
 Steklov Inst. Math.}} \textbf{2} (2016), 139--153, \href{https://arxiv.org/abs/1603.01706}{arXiv:1603.01706}.

\end{thebibliography}\LastPageEnding
